\section{证伪 ETH 的其他推论}\label{app:consequences}
定理~\ref{thm:main} 给出一个电路下界，以及求逆常数输出
局部性布尔映射的亚指数时间算法。固定密度情形给出
推论~\ref{cor:quantitative-circuits} 的定量电路下界。
两个电路论证都只用我们 $3$-SAT 算法的判定版本；
精确原像计数使用其完整的计数保证。

\subsection{$\mathrm E^{\mathrm{NP}}$ 的电路下界}\label{subsec:circuit-lower-bounds}
对下文的定性界，记 $\mathrm{SIZE}(O(n))$
为拥有 $O(n)$ 个有界扇入 AND、OR、NOT 门的非一致电路的
语言类。

\begin{corollary}\label{cor:consequence-circuits}
我们有
\[
 \mathrm E^{\mathrm{NP}}\not\subseteq\mathrm{SIZE}(O(n)).
\]
\end{corollary}
\begin{proof}
Williams 证明：若 ETH 为假，则
$\mathrm E^{\mathrm{NP}}$ 没有线性规模电路
~\cite[Theorem~6.1]{Wil13}。把该蕴含应用于
定理~\ref{thm:main} 提供的确定性亚指数时间 $3$-SAT
算法。
\end{proof}

这一联系是从更快的可满足性算法导出电路下界的一个实例。
为编码有 $s$ 个门、$n$ 个输入的布尔电路，
为每个门引入一个变量，用其输入强制其取值，
并添加一个要求输出位等于 $1$ 的单元子句。
所得 $3$-CNF 有 $O(n+s)$ 个变量与子句，
且恰在电路输出为一时可满足。
因此我们的定理在 $2^{o(n)}$ 时间内判定线性规模电路的
可满足性。Williams 的论证把这样的加速与对见证的假想小
电路结合，与某个非确定性时间层级矛盾。
该推论给出 $\mathrm E^{\mathrm{NP}}$
中没有线性规模电路族的语言；
它不断言 NP 的超多项式电路下界。

\input{quantitative-circuits-zh}

\subsection{求逆常数输出局部性的布尔映射}\label{subsec:local-inversion}
固定常数 $d$。考虑布尔映射
$G:\{0,1\}^n\to\{0,1\}^m$，其第 $i$ 个输出位形如
\[
 G(\vec{x})_i=g_i(\vec{x}|_{S_i}),
 \qquad S_i\subseteq[n],\quad |S_i|\le d.
\]
$G$ 的输入描述列出集合 $S_i$ 与布尔函数 $g_i$
的真值表。于是每个输出只检查常数个输入位，
尽管集合 $S_i$ 可以任意重叠。
给定目标 $\vec{b}\in\{0,1\}^m$，
求逆即要求任意满足 $G(\vec{x})=\vec{b}$
的 $\vec{x}$。

\begin{corollary}\label{cor:consequence-local-inversion}
对每个固定的 $d$ 与 $n\ge20$，
数 $|G^{-1}(\vec{b})|$ 可以被精确计算，
且当原像存在时可以找到其一，
确定性时间为
\[
 (n+m)^{O_d(1)}2^{O_d(\rho(n))}
\]
，空间为 $(n+m)^{O_d(1)}$。
\end{corollary}
\begin{proof}
对每个 $i$ 与每个满足 $g_i(\vec{\alpha})\ne b_i$
的局部赋值 $\vec{\alpha}\in\{0,1\}^{S_i}$，
加入子句
\[
 \bigvee_{j\in S_i}\ell_{j,\alpha_j},
 \qquad
 \ell_{j,0}=X_j,\quad \ell_{j,1}=\neg X_j.
\]
该子句恰禁止局部赋值 $\vec{\alpha}$。
若 $S_i$ 为空，其约束是常值的，可直接检验。
剩余子句的合取是原 $n$ 个声明变量上的 CNF，
宽度至多 $d$、子句至多 $2^d m$ 个。
其可满足赋值恰为 $G^{-1}(\vec{b})$；
不需要辅助变量或重数修正。
以 $k=\max\{2,d\}$ 应用定理~\ref{thm:main}。
通常的逐变量自归约用至多 $n$ 次进一步计数查询找到一个
原像，这只改变多项式因子。
每次查询可在下一次之前丢弃，
保持多项式空间。
\end{proof}

条件计数还允许在原像集非空时采样一个完全均匀的原像，
期望时间界相同（差多项式因子）~\cite{JVV86}。
特别地，该推论适用于 Goldreich 候选单向函数的常数
局部性版本，其在输入位的重叠子集上求值局部谓词
~\cite{Gol00}。对多项式个输出位，
它给出对像中每个输出的 $2^{o(n)}$ 时间求逆算法，
排除这些族的 $2^{\Omega(n)}$ 求逆难度。
这并不排除对多项式时间算法的单向性。

\input{quantitative-reduction-zh}
\input{quantitative-contraction-zh}

